\section{Details of Non-Implications}
\label{sec:cex-extra}

\subsection{Trivial Examples}
\label{sec:cex-extra:trivial}

\begin{example}[single item]
\label{cex:single-item}
Consider a fair division instance with $n$ agents and one item
(which is either a good to everyone or a chore to everyone).
Then every allocation is EFX, EF1, APS, MMS, PROPx, GAPS, GMMS, PAPS, PMMS,
EEFX, EEF1, MXS, M1S, PROPm, and PROP1,
but not EF or PROP or GPROP or PPROP or EEF or MEFS.
\end{example}

\begin{lemma}
\label{cex:share-vs-envy-goods}
Consider a fair division instance $([n], [m], (v_i)_{i=1}^n, \eqEnt)$
with $n \ge 3$, $m = 2n-1$, and identical additive valuations,
where each item has value 1 to each agent.
Let $A$ be an allocation where agent $n$ gets $n$ goods,
and all other agents get 1 good each.
Then this allocation is APS+MMS+EEFX+PROPx, but not EF1.
\end{lemma}
\begin{proof}
By \cref{thm:impl:aps-to-pess,thm:impl:prop-to-aps},
$1 = \MMS_i \le \APS_i \le v([m])/n = 2 - 1/n$.
Hence, $APS_i = 1$, since the APS is the value of some bundle.
Hence, $A$ is APS+MMS+PROPx.
Agent $n$'s EEFX-certificate for $A$ is $A$ itself.
For $i \neq n$, agent $n$'s EEFX-certificate is $B$,
where $|B_i| = 1$ and $|B_j| = 2$ for $j \neq i$.
%
$A$ is not EF1 because the first $n-1$ agents EF1-envy agent $n$.
\end{proof}

\begin{lemma}
\label{cex:share-vs-envy-chores}
Consider a fair division instance $([n], [m], (v_i)_{i=1}^n, w)$
with $n \ge 3$, equal entitlements, identical additive disutilities,
and $m = n+1$ chores, each of disutility 1.
Let $A$ be an allocation where agents 1 and 2 get 2 chores each,
agent $n$ gets 0 chores, and the remaining agents get 1 chore each.
Then this allocation is APS+MMS+EEFX+PROPx, but not EF1.
\end{lemma}
\begin{proof}
$-2 = \MMS_i \le \APS_i \le v([m])/n = - 1 - 1/n$.
Hence, $APS_i = -2$, since the APS is the value of some bundle.
Hence, $A$ is APS+MMS+PROPx.

Agents $[n] \setminus [2]$ do not EFX-envy anyone in $A$.
For $i \in \{1, 2\}$, agent $i$'s epistemic-EFX-certificate can be obtained by
transferring a chore from agent $3-i$ to agent $n$.
Hence, $A$ is epistemic EFX.

$A$ is not EF1 because agents 1 and 2 EF1-envy agent $n$.
\end{proof}

\subsection{From EEF, MEFS, PROP}
\label{sec:cex-extra:from-eef-mefs-prop}

\begin{lemma}[EEF $\nfimplies$ EF1]
\label{cex:eef-not-ef1}
Let $0 \le 2a < b$. Let $f_1, f_2, f_3: 2^{[12]} \to \mathbb{R}_{>0}$ be additive sets functions:

\begin{tabular}{c|cccc|cccc|cccc}
& 1 & 2 & 3 & 4 & 5 & 6 & 7 & 8 & 9 & 10 & 11 & 12
\\ \hline $f_1$ & $a$ & $a$ & $b$ & $b$ & $b$ & $b$ & $b$ & $b$ & $a$ & $a$ & $a$ & $a$
\\ $f_2$ & $a$ & $a$ & $a$ & $a$ & $a$ & $a$ & $b$ & $b$ & $b$ & $b$ & $b$ & $b$
\\ $f_3$ & $b$ & $b$ & $b$ & $b$ & $a$ & $a$ & $a$ & $a$ & $a$ & $a$ & $b$ & $b$
\end{tabular}

Let $t \in \{-1, 1\}$ and let $\Ical \defeq ([3], [12], (v_i)_{i=1}^3, \eqEnt)$
be a fair division instance where $v_i \defeq tf_i$ for all $i \in [3]$.
Then allocation
$A \defeq ([4], [8] \setminus [4], [12] \setminus [8])$ is EEF+PROP but not EF1.
\end{lemma}
\begin{proof}
For $t = 1$, agent 1 EF1-envies agent 2 in $A$,
and for $t = -1$, agent 1 EF1-envies agent 3 in $A$.
$B = ([4], \{5, 6, 9, 10\}, \{7, 8, 11, 12\})$ is agent 1's EEF-certificate.
A similar argument holds for agents 2 and 3 too.
\end{proof}

\begin{example}[PROP $\nfimplies$ MEFS]
\label{cex:prop-not-mefs-goods}
Consider a fair division instance with 3 equally-entitled agents
having additive valuations over 3 goods:

\begin{tabular}{c|ccc}
& 1 & 2 & 3
\\ \hline $v_1$ & 10 & 20 & 30
\\ $v_2$ & 20 & 10 & 30
\\ $v_3$ & 10 & 20 & 30
\end{tabular}

Then the allocation $(\{2\}, \{1\}, \{3\})$ is PROP, but no allocation is MEFS
(every agent's minimum EF share is 30).
\end{example}

\begin{example}[PROP $\nfimplies$ MEFS]
\label{cex:prop-not-mefs-chores}
Consider a fair division instance with 3 equally-entitled agents
having additive disutilities over 3 chores:

\begin{tabular}{c|ccc}
& 1 & 2 & 3
\\ \hline $-v_1$ & 30 & 20 & 10
\\ $-v_2$ & 20 & 30 & 10
\\ $-v_3$ & 30 & 20 & 10
\end{tabular}

Then the allocation $(\{2\}, \{1\}, \{3\})$ is PROP, but no allocation is MEFS
(every agent's minimum EF share is $-10$).
\end{example}

\begin{lemma}[MEFS $\nfimplies$ EEF]
\label{cex:mefs-not-eef-goods}
Consider a fair division instance with 3 equally-entitled agents
having additive valuations over 6 goods:

\begin{tabular}{c|cccccc}
& 1 & 2 & 3 & 4 & 5 & 6
\\ \hline $v_1$ & 20 & 20 & 20 & 10 & 10 & 10
\\ $v_2$, $v_3$ & 20 & 10 & 10 &  1 &  1 &  1
\end{tabular}

Then the allocation $A \defeq (\{4, 5, 6\}, \{1\}, \{2, 3\})$ is MEFS, but no allocation is epistemic EF.
\end{lemma}
\begin{proof}
Agents 2 and 3 are envy-free in $A$.
Agent 1 has $B \defeq (\{1, 4\}, \{2, 5\}, \{3, 6\})$ as her MEFS-certificate for $A$.
Hence, $A$ is MEFS.

Suppose an epistemic EF allocation $X$ exists.
Let $Y^{(i)}$ be each agent $i$'s epistemic-EF-certificate.
For agent 2 to be envy-free in $Y^{(2)}$,
we require $Y^{(2)}_2 \supseteq \{1\}$ or $Y^{(2)}_2 \supseteq \{2, 3\}$.
Similarly, $Y^{(3)}_3 \supseteq \{1\}$ or $Y^{(3)}_3 \supseteq \{2, 3\}$.
Since $Y^{(i)}_i = X_i$ for all $i$, we get $X_2 \cup X_3 \supseteq \{1, 2, 3\}$.
Hence, $X_1 \subseteq \{4, 5, 6\}$.
But then no epistemic-EF-certificate exists for agent 1 for $X$,
contradicting our assumption that $X$ is epistemic EF.
Hence, no epistemic EF allocation exists.
\end{proof}

\begin{example}[MEFS $\nfimplies$ EEF]
\label{cex:mefs-not-eef-chores}
Consider a fair division instance with 3 equally-entitled agents
having additive disutilities over 6 chores:

\begin{tabular}{c|cccccc}
& 1 & 2 & 3 & 4 & 5 & 6
\\ \hline $-v_1$ & 20 & 20 & 20 & 10 & 10 & 10
\\ $-v_2$, $-v_3$ & 20 & 10 & 10 & 10 & 10 & 10
\end{tabular}

Then the allocation $A \defeq (\{4, 5, 6\}, \{1\}, \{2, 3\})$ is MEFS
(agents 2 and 3 are EF, agent 1's MEFS-certificate is $(\{1, 4\}, \{2, 5\}, \{3, 6\})$).
Agent 1 is not epistemic-EF-satisfied by $A$.
\end{example}

\begin{lemma}[MEFS $\nfimplies$ EEF1]
\label{cex:mefs-not-eef1-chores}
Consider a fair division instance with 3 equally-entitled agents
having additive disutilities over 12 chores.
$v_1(1) = v(2) = v(3) = 70$ and $v(c) = 10$ for all $c \in [12] \setminus [3]$.
Agents 2 and 3 have disutility 10 for each chore.
Then $A \defeq ([12] \setminus [3], [2], \{3\})$ is a MEFS+PROP allocation
where agent 1 is not EEF1-satisfied.
\end{lemma}
\begin{proof}
$\PROP_1 = -100$ and $\PROP_2 = \PROP_3 = -40$.
$\MEFS_1 \le -100$ because of the allocation $(\{1, 4, 5, 6\}, \{2, 7, 8, 9\}, \{3, 10, 11, 12\})$.
$\MEFS_i \le -40$ for $i \in \{2, 3\}$ because of the allocation
$([4], [8] \setminus [4], [12] \setminus [8])$.
Agent 1 has disutility $90$ in $A$, so $A$ is MEFS-fair and PROP-fair to agent 1.
Agents 2 and 3 have disutility at most $20$ in $A$, so $A$ is MEFS-fair and PROP-fair to them.

Agent 1 is not EEF1-satisfied by $A$, since in any EEF1-certificate $B$,
some agent $j \in \{2, 3\}$ receives at most one chore of value $70$,
and agent 1 would EF1-envy $j$.
\end{proof}

\subsection{Two Equally-Entitled Agents}
\label{sec:cex-extra:2-eqEnt}

\begin{example}[EFX $\nfimplies$ MMS]
\label{cex:efx-not-mms}
Let $t \in \{-1, 1\}$.
Consider a fair division instance with 2 equally-entitled agents having
an identical additive valuation function $v$ over 5 items.
$v(1) = v(2) = 3t$ and $v(3) = v(4) = v(5) = 2t$.
Then allocation $A \defeq (\{1, 3\}, \{2, 4, 5\})$ is EFX.
The MMS is $6t$, since $P = (\{3t, 3t\}, \{2t, 2t, 2t\})$ is an MMS partition.
But in $A$, some agent doesn't get her MMS.
\end{example}

\begin{example}[EF1 $\nfimplies$ PROPX or MXS]
\label{cex:ef1-not-propx-mxs}
Let $t \in \{-1, 1\}$.
Consider a fair division instance with 2 equally-entitled agents
having an identical additive valuation function $v$ over 5 items,
where $v(1) = v(2) = 4t$ and $v(3) = v(4) = v(5) = t$.
Then allocation $A \defeq (\{1\}, [5] \setminus \{1\})$ is EF1 but not PROPx and not MXS.
\end{example}

\begin{lemma}[PROPx $\nfimplies$ M1S]
\label{cex:propx-not-m1s}
Let $t \in \{-1, 1\}$ and $0 < \eps < 1/2$.
Consider a fair division instance with 2 equally-entitled agents having
an identical additive valuation function $v$ over 4 items.
Let $v(4) = (1+2\eps)t$ and $v(j) = t$ for $j \in [3]$.
Then allocation $A \defeq (\{4\}, [3])$ is PROPx but not M1S.
\end{lemma}
\begin{proof}
$v([m])/2 = (2+\eps)t$, so $A$ is PROPx.
%
For $t = 1$, in any allocation $B$ where agent 1 doesn't EF1-envy agent 2, she must have at least 2 goods.
But $v(A_1) = 1+2\eps$, so agent 1 doesn't have an M1S-certificate for $A$. Hence, $A$ is not M1S.
%
For $t = -1$, in any allocation $B$ where agent 2 doesn't EF1-envy agent 1, she must have at most 2 chores.
But $v(A_2) = -3$, so agent 1 doesn't have an M1S-certificate for $A$. Hence, $A$ is not M1S.
\end{proof}

\begin{example}[MXS $\nfimplies$ PROPx for $n=2$, \cite{caragiannis2022existence}]
\label{cex:mxs-not-propx-n2}
Let $t \in \{-1, 1\}$.
Consider a fair division instance with 2 equally-entitled agents
having identical additive valuations over 7 items:
the first 2 items of value $4t$ and the last 5 items of value $t$.
Then the allocation $A = (\{1, 3\}, \{2, 4, 5, 6, 7\})$ is not PROPx or EFX,
but it is MXS because the agents have $([7] \setminus [2], [2])$
and $([2], [7] \setminus [2])$ as their MXS-certificates for $A$.
\end{example}

\begin{lemma}[M1S $\nfimplies$ PROP1]
\label{cex:m1s-not-prop1}
Consider a fair division instance with 2 equally-entitled agents
having an identical additive valuation function $v$ over 9 items.
Let $t \in \{-1, 1\}$ and $v(9) = 4t$ and $v(j) = t$ for $j \in [8]$.
Then allocation $A \defeq (\{9\}, [8])$ is M1S but not PROP1.
\end{lemma}
\begin{proof}
$v([9])/2 = 6t$. Let $B \defeq ([4], [9] \setminus [4])$.
%
For $t = 1$ (goods), $B$ is agent 1's M1S-certificate for $A$,
but agent 1 is not PROP1-satisfied by $A$.
%
For $t = -1$ (chores), $B$ is agent 2's M1S-certificate for $A$,
but agent 2 is not PROP1-satisfied by $A$.
\end{proof}

\subsection{Three Equally-Entitled Agents}
\label{sec:cex-extra:3-eqEnt}

\begin{example}[GAPS $\nfimplies$ PROPx]
\label{cex:gaps-not-propx}
Consider a fair division instance with 3 equally-entitled agents
having identical additive or unit-demand valuations. There are 2 goods of values 50 and 10.
In every allocation, some agent doesn't get any good, and that agent is not PROPx-satisfied.
The allocation where the first agent gets the good of value 50
and the second agent gets the good of value 10 is a groupwise APS allocation
(set the price of the goods to $1.1$ and $0.9$).
\end{example}

\begin{lemma}[APS $>$ MMS]
\label{thm:aps-gt-mms}
Let $t \in \{-1, 1\}$.
Consider a fair division instance with 3 equally-entitled agents
having identical additive valuations over 15 items. The items' values are
$65t$, $31t$, $31t$, $31t$, $23t$, $23t$, $23t$, $17t$, $11t$, $7t$, $7t$, $7t$, $5t$, $5t$, $5t$.
Then the AnyPrice share is at least $97t$, the proportional share is $97t$,
and the maximin share is less than $97t$.
\end{lemma}
\begin{proof}
For $t = 1$, this follows from Lemma C.1 of \cite{babaioff2023fair}.
For $t = -1$, a similar argument tells us that the AnyPrice share is at least $-97$.
If the maximin share is at least $-97$, then there must exist a partition $P$ of the chores
where each bundle has disutility 97. But then $P$ would prove that the maximin share
in the corresponding goods instance is at least 97, which is a contradiction.
Hence, for $t = -1$, the maximin share is less than $-97$.
\end{proof}

\begin{example}[GMMS $\nfimplies$ APS]
\label{cex:gmms-not-aps}
For the fair division instance in \cref{thm:aps-gt-mms},
the leximin allocation is GMMS (since on restricting to any subset of agents,
the resulting allocation is still leximin, and is therefore MMS).
However, no APS allocation exists, because APS $>$ MMS,
and the minimum value across all bundles is at most the MMS.
\end{example}

\begin{example}[PMMS $\nfimplies$ MMS, Example 4.4 of \cite{caragiannis2019unreasonable}]
\label{cex:pmms-not-mms}
Let $t \in \{-1, 1\}$.
Consider a fair division instance with 3 equally-entitled agents
having an identical valuation function $v$ over 7 items where
$v(1) = 6t$, $v(2) = 4t$, $v(3) = v(4) = 3t$, $v(5) = v(6) = 2t$, $v(7) = t$.
Each agent's maximin share is $7t$ ($(\{1, 7\}, \{2, 3\}, \{4, 5, 6\})$ is a maximin partition).
Allocation $(\{1\}, \{3, 4, 5\}, \{2, 6, 7\})$ is PMMS but not MMS.
\end{example}

\begin{example}[APS $\nfimplies$ PROPm]
\label{cex:aps-not-propm}
Consider a fair division instance with 3 equally-entitled agents
having an identical additive valuation function $v$ over 6 goods:
$v(1) = 60$, $v(2) = 30$, and $v(3) = v(4) = v(5) = v(6) = 10$.
The allocation $A \defeq (\{2\}, \{3, 4, 5\}, \{1, 6\})$ is APS+MMS, since the MMS is 30,
and the APS is at most 30 because of the price vector $(4, 3, 1, 1, 1, 1)$.
However, $A$ is not PROPm-fair to agent 1, because the proportional share is $130/4 > 40$.
\end{example}

\begin{example}[APS $\nfimplies$ PROP1]
\label{cex:aps-not-prop1-chores}
Consider a fair division instance with 3 equally-entitled agents
having an identical additive valuation function $v$ over 6 chores:
the first chore has disutility 18 (large chore) and the remaining chores have disutility 3 each (small chores).
Then $X \defeq ([6] \setminus \{1\}, \{1\}, \emptyset)$ is MMS+APS,
since the MMS is $-18$, and the APS is at most $-18$ due to the price vector $(1, 0, 0, 0, 0, 0)$.
$X$ is not PROP1-fair to agent 1, since the proportional share is $-11$,
and agent 1's disutility in $X$ after removing any chore is $12$.
$X$ is not EEF1-fair to agent 1 because even after redistributing chores among the remaining agents,
someone will always have no chores.
\end{example}

\begin{lemma}[PROPm doesn't exist for mixed manna]
\label{cex:propm-mixed-manna}
Consider a fair division instance $([3], [6], (v_i)_{i=1}^3, \eqEnt)$
where agents have identical additive valuations, and the items have values
$(-3, -3, -3, -3, -3, 3\eps)$, where $0 < \eps < 1/2$.
Then there exists an EFX+GMMS+GAPS allocation but no PROPm allocation.
\end{lemma}
\begin{proof}
The proportional share is $v([m])/3 = -5 + \eps$.
\WLoG, assume agent 1 receives the most number of chores,
and agent 3 receives the least number of chores.
Then agent 1 has at least 2 chores, and agent 3 has at most 1 chore.

\textbf{Case 1}: agent 1 receives at least 3 chores.
\\ Then even after removing one of her chores, and even if she receives the good,
her value for her bundle is at most $-6 + 3\eps < -5 + \eps$.
Hence, she is not PROPm-satisfied.

\textbf{Case 2}: agent 1 receives 2 chores.
\\ Then agent 2 also receives 2 chores, and agent 3 receives 1 chore.
Assume without loss of generality that agent 2 receives the good.
Then this allocation is EFX and groupwise MMS.
On setting the price of each chore to $-3$ and the price of the good to 3,
we get that the allocation is groupwise APS.
However, if agent 1 adds the good to her bundle, her value becomes $-6 + 3\eps < -5 + \eps$.
Hence, she is not PROPm-satisfied.
\end{proof}

\begin{example}[PROPm $\nfimplies$ PROPavg]
\label{cex:propm-not-propavg}
Consider a fair division instance with three equally-entitled agents
having an identical additive valuation function $v$ over three goods.
Let $v(1) = v(2) = 60$ and $v(3) = 30$.
Then allocation $A \defeq (\{3\}, \{1, 2\}, \emptyset)$ is PROPm,
but agent 3 is not PROPavg-satisfied by $A$.
\end{example}

\subsection{Unequal Entitlements}
\label{sec:cex-extra:uneqEnt}

\begin{lemma}[PROP1+M1S is infeasible]
\label{cex:prop1-plus-m1s-ue}
Consider a fair division instance $\Ical \defeq ([3], [7], (v_i)_{i=1}^3, w)$,
where the entitlement vector is $w \defeq (7/12, 5/24, 5/24)$,
the agents have identical additive valuations, and each good has value 1. Then
\begin{tightenum}
\item $\APS_1 = 4$ and $\APS_2 = \APS_3 = 1$.
\item $X$ is APS $\iff$ $X$ is groupwise-APS (GAPS) $\iff$ $X$ is PROP1.
\item $\WMMS_1 = 3$ and $\WMMS_2 = \WMMS_3 = 15/14$.
\item $X$ is WMMS $\iff$ $X$ is groupwise-WMMS (GWMMS) $\iff$ $X$ is EFX $\iff$ $X$ is M1S.
\end{tightenum}
Therefore,
\begin{tightenum}
\item M1S+PROP1 is infeasible for this instance.
\item GWMMS+EFX doesn't imply PROP1.
\item GAPS doesn't imply M1S.
\end{tightenum}
\end{lemma}
\begin{proof}
By \cref{thm:impl:tribool:aps}, $\APS_1 = \floor{\frac{7 \times 7}{12}} = 4$
and $\APS_2 = \APS_3 = \floor{\frac{5 \times 7}{24}} = 1$.
By \cref{thm:impl:tribool:aps,thm:impl:tribool:prop1},
an allocation is APS iff it is PROP1.

Any GAPS allocation is also APS by definition.
We will now show that any APS allocation is also GAPS.
Formally, let $A$ be an APS allocation for $\Ical$.
The cardinality vector of $A$, i.e., $c \defeq (|A_1|, |A_2|, |A_3|)$,
can have three possible values: $(5, 1, 1)$, $(4, 2, 1)$, $(4, 1, 2)$.
For every possible value of $c$ and $S \subseteq [3]$,
we show that $(\Icalhat, \Ahat) \defeq \restrict(\Ical, A, S)$ is APS
(c.f.~\cref{defn:restricting}).
\begin{tightenum}
\item $c = (5, 1, 1)$ and $S = \{1, 2\}$:
    $\Icalhat$ has 6 goods and entitlement vector $(14/19, 5/19)$.
    $\APS_1 = \floor{\frac{14 \times 6}{19}} = 4$ and $\APS_2 = \floor{\frac{5 \times 6}{19}} = 1$.
    Hence, $\Ahat$ is APS for $\Icalhat$.
\item $c = (5, 1, 1)$ and $S = \{1, 3\}$:
    Similar to the $S = \{1, 2\}$ case.
\item $c = (5, 1, 1)$ and $S = \{2, 3\}$:
    $\Icalhat$ has 2 goods and entitlement vector $(1/2, 1/2)$.
    $\APS_3 = \APS_3 = 1$, so $\Ahat$ is APS for $\Icalhat$.
\item $c = (4, 2, 1)$ and $S = \{1, 2\}$:
    $\Icalhat$ has 6 goods and entitlement vector $(14/19, 5/19)$.
    $\APS_1 = \floor{\frac{14 \times 6}{19}} = 4$ and $\APS_2 = \floor{\frac{5 \times 6}{19}} = 1$.
    Hence, $\Ahat$ is APS for $\Icalhat$.
\item $c = (4, 2, 1)$ and $S = \{1, 3\}$:
    $\Icalhat$ has 5 goods and entitlement vector $(14/19, 5/19)$.
    $\APS_1 = \floor{\frac{14 \times 5}{19}} = 3$ and $\APS_2 = \floor{\frac{5 \times 5}{19}} = 1$.
    Hence, $\Ahat$ is APS for $\Icalhat$.
\item $c = (4, 2, 1)$ and $S = \{2, 3\}$:
    $\Icalhat$ has 3 goods and entitlement vector $(1/2, 1/2)$.
    $\APS_2 = \APS_3 = \floor{\frac{1 \times 3}{2}} = 1$.
    Hence, $\Ahat$ is APS for $\Icalhat$.
\item $c = (4, 1, 2)$:
    Similar to the $c = (4, 2, 1)$ case.
\end{tightenum}
Hence, any APS allocation for $\Ical$ is also GAPS.

For any allocation $X$, define
\[ f(X) \defeq \min_{j=1}^3 \frac{|X_j|}{w_j}. \]
Then $\WMMS_i = w_i\max_X f(X)$ for all $i \in [3]$.
If $|X_1| \le 2$, then $f(X) \le |X_1|/w_1 \le 24/7 = 3 + 3/7$.
If $|X_2| \le 1$ or $|X_3| \le 1$, $f(X) \le 24/5 = 4 + 4/5$.
Otherwise, $|X_1| = 3$ and $|X_2| = |X_3| = 2$,
so $f(X) = \min(3 \times 12/7, 2 \times 24/5) = 36/7 = 5 + 1/6$.
Hence, $\max_X f(X) = 36/7$, so $\WMMS_1 = 3$ and $\WMMS_2 = \WMMS_3 = 15/14$.
So, an allocation is WMMS iff it has cardinality vector $(3, 2, 2)$.

Any GWMMS allocation is also WMMS by definition. We now prove the converse.
For every possible value of $S \subseteq [3]$,
we show that $(\Icalhat, \Ahat) \defeq \restrict(\Ical, A, S)$ is WMMS
(c.f.~\cref{defn:restricting}).
\begin{tightenum}
\item $S = \{1, 2\}$:
    $\Icalhat$ has 5 goods and entitlement vector $(14/19, 5/19)$.
    If $|X_1| \le 2$, then $f(X) \le |X_1|/w_1 \le 38/14 = 2 + 10/14$.
    If $|X_2| \le 1$, then $f(X) \le |X_2|/w_2 \le 19/5 = 3 + 10/14$.
    Otherwise, $|X_1| = 3$ and $|X_2| = 2$, so $f(X) = \min(3/w_1, 2/w_2) = \min(57/14, 38/5) = 57/14 = 4 + 1/14$.
    Hence, $\WMMS_1 = 3$ and $\WMMS_2 = \WMMS_3 = 57/14 \times 5/19 = 1 + 19/266$.
    Hence, $\Ahat$ is WMMS for $\Icalhat$.
\item $S = \{1, 3\}$:
    Similar to the $S = \{1, 2\}$ case.
\item $S = \{2, 3\}$:
    $\Icalhat$ has 4 goods and entitlement vector $(1/2, 1/2)$.
    Then $\WMMS_1 = \WMMS_2 = 2$.
    Hence, $\Ahat$ is WMMS for $\Icalhat$.
\end{tightenum}
Hence, any WMMS allocation for $\Ical$ is also GWMMS.

Any GWMMS allocation is EFX by \cref{thm:impl:mms-to-efx-n2},
and any EFX allocation is M1S by \cref{thm:impl:efx-to-ef1}.
We will now show that any M1S allocation is WMMS.

Let $X$ be an M1S allocation.
Let $A$ be agent 1's M1S certificate.
If $|A_1| \le 2$, then $|A_j| \ge 3$ for some $j \in \{2, 3\}$.
Since agent 1 has higher entitlement, she would EF1-envy agent $j$,
which is a contradiction. Hence, $|X_1| \ge |A_1| \ge 3$.

Let $B$ be agent 2's M1S certificate.
Suppose $|B_2| \le 1$. Since $2$ doesn't EF1-envy $3$, we get $|B_3| \le 2$.
Since $2$ doesn't EF1-envy $1$, we get
\[ \frac{|B_1|-1}{w_1} \le \frac{|B_2|}{w_2} \iff |B_1| \le 1 + \frac{w_1}{w_2} = 3 + \frac{4}{5}. \]
Hence, $|B_1| + |B_2| + |B_3| \le 3 + 1 + 2 = 6$, which is a contradiction.
Hence, $|X_2| \ge |B_2| \ge 2$.
Similarly, we can prove that $|X_3| \ge 2$.

Hence, $|X_i| \ge \WMMS_i$ for all $i$, so $X$ is WMMS.
This proves that any M1S allocation is WMMS.
\end{proof}

\cite{chakraborty2021weighted} also proves that EF1+PROP1 allocations may not exist for unequal entitlements,
We use a different counterexample in \cref{cex:prop1-plus-m1s-ue},
which allows us to also prove other non-implications.

\begin{example}[PROP1 $\nfimplies$ M1S]
\label{cex:prop1-not-m1s-n2}
Consider a fair division instance with 2 agents having identical additive valuations.
Let $t \in \{-1, 1\}$. Let there be 2 items, each of value $t$.
Let the entitlement vector be $(2/3, 1/3)$.
Let $A$ be an allocation where the first agent gets both items.
Then $A$ is PROP1 but not M1S.
\end{example}

\begin{lemma}[GAPS $\nfimplies$ PROPx for $n=2$]
\label{cex:gaps-not-propx-n2}
Let $t \in \{-1, 1\}$. Consider a fair division instance $([2], [12], (v, v), (0.4, 0.6))$,
where $v(1) = v(2) = 10t$ and $v(j) = t$ for all $j \in [12] \setminus [2]$.
Then allocation $A = (\{1\}, [12] \setminus \{1\})$ is APS but not PROPx.
\end{lemma}
\begin{proof}
Agents 1 and 2 have proportional shares $12t$ and $18t$, respectively.

For $t = 1$, agent 1's APS is at most $10$
(price items 1 and 2 at $0.45$ each and the remaining items at $0.01$ each),
and agent 2's APS is at most $18$ (price each item at its value).
Hence, $A$ is APS but not PROPx-fair to agent 1.

For $t = -1$, agent 2's APS is at most $-20$
(price items 1 and 2 at $-0.45$ each and the remaining items at $-0.01$ each),
and agent 1's APS is at most $-12$ (price each item at its value).
Hence, $A$ is APS but not PROPx-fair to agent 2.
\end{proof}

\subsection{Supermodular Valuations}
\label{sec:cex-extra:supmod}

\begin{lemma}[EF+GAPS $\nfimplies$ PROP1]
\label{cex:ef-not-prop-supmod}
Let $0 \le \eps < 1/27$ and $k \in \mathbb{N}$. For any $x \ge 0$, let
\[ f(x) \defeq x\eps + \begin{cases}
0 & \textrm{ if } x \le k
\\ x-k & \textrm{ if } x > k
\end{cases}. \]
Consider a fair division instance with $n$ equally-entitled agents
having an identical valuation function $v$ over $4n$ goods,
where $v(X) = f(|X|)$ for all $X \subseteq [4n]$. Then $v$ is supermodular.

Let $A$ be the allocation where each agent gets 4 goods. Then $A$ is EF+GAPS+GMMS.
If $k = 5$, then no PROP1 or PROPm allocation exists.
If $n \ge 3$ and $k = 8$, then $A$ is PPROP but no PROP1 or PROPx allocation exists.
Additionally if $\eps = 0$, then $A$ is PROPavg.
\end{lemma}
\begin{proof}
For any $X \subseteq [4n]$ and $g \in [4n] \setminus X$, we have
\[ v(g \mid X) = \begin{cases}
\eps & \textrm{ if } |X| \le k
\\ 1 + \eps & \textrm{ if } |X| \ge k
\end{cases}. \]
Since $v(g \mid X)$ is non-decreasing in $|X|$, $v$ is supermodular.

Since all agents have the same bundle, $A$ is EF and GMMS.
Set the price of each good to 1 to get that $A$ is GAPS.

The proportional share is $4-k/n+4\eps$.
For $k = 5$, the proportional share is at least $3/2$,
but $f(5) = 5\eps$ and $f(6) = 1+6\eps < 3/2$.
In any allocation, some agent will have at most 4 goods,
and that agent would not be PROP1 or PROPm satisfied.

If $n \ge 3$ and $k = 8$, then the proportional share is at least $4/3$.
Since $f(8) = 8\eps$ and $f(4) = 4\eps$, $A$ is pairwise PROP.
In any allocation, some agent gets at most 4 goods,
and $f(5) = 5\eps < 4/3$, so that agent is not PROP1-satisfied,
and $f(9) = 1 + 9\eps < 4/3$, so that agent is not PROPx-satisfied.
Additionally, if $\eps = 0$, for any two agents $i$ and $j$,
we have $v(A_j \mid A_i) = f(8) - f(4) = 0$. Hence, $A$ is PROPavg.
\end{proof}

\subsection{Unit-Demand Valuations}
\label{sec:cex-extra:unit-demand}

It is known that unit-demand functions are submodular and cancelable.
We first show how to perturb a unit-demand function such that
submodularity and cancelability remain preserved and all marginals become positive.

\begin{lemma}
\label{thm:ud-perturb}
Let $u: 2^{[m]} \to \mathbb{R}_{\ge 0}$ be a function and $\eps \in \mathbb{R}_{\ge 0}$.
For any $X \subseteq [m]$, define $\uhat(X) \defeq u(X) + |X|\eps$.
If $u$ is submodular, then $\uhat$ is also submodular.
If $m \le 3$ and $u$ is cancelable, then $\uhat$ is also cancelable.
\end{lemma}
\begin{proof}
If $u$ is submodular, then for any $S, T \subseteq [m]$, we get
\[ \uhat(S) + \uhat(T) - \uhat(S \cup T) - \uhat(S \cap T)
= u(S) + u(T) - u(S \cup T) - u(S \cap T) \ge 0. \]
Hence, $\uhat$ is also submodular.

Let $u$ be cancelable.
For every $T \subseteq [m]$ and $S_1, S_2 \subseteq [m] \setminus T$,
if $\uhat(S_1 \cup T) > \uhat(S_2 \cup T)$,
we need to show that $\uhat(S_1) > \uhat(S_2)$
to prove that $\uhat$ is cancelable.
This is trivial when $T = \emptyset$ or $\eps = 0$.

If $|S_1| = |S_2|$, then $u(S_1 \cup T) > u(S_2 \cup T)$.
Since $u$ is cancelable, we get $u(S_1) > u(S_2)$, which implies $\uhat(S_1) > \uhat(S_2)$.

Let $\eps > 0$. If $S_1 = \emptyset$, we get
$\uhat(S_1 \cup T) > \uhat(S_2 \cup T) \implies u(S_2 \mid T) < - |S_2|\eps$,
which is a contradition since $u$ is non-negative. Hence, $S_1 \neq \emptyset$.
If $S_2 = \emptyset$, we get $\uhat(S_1) \ge |S_1|\eps > 0 = \uhat(S_2)$.
When $m \le 3$, we have considered all cases.
Hence, if $m \le 3$ and $u$ is cancelable, then $\uhat$ is also cancelable.
\end{proof}

\begin{lemma}[PROP $\nfimplies$ M1S]
\label{cex:prop-not-m1s-unit-demand}
Consider a unit-demand function $v$ over 3 goods such that $v(1) = v(2) = 4$ and $v(3) = 3$.
Let $0 \le \eps < 1/6$ and $\vhat(X) \defeq v(X) + 2|X|\eps$ for all $X \subseteq [3]$.
Consider a fair division instance with 3 goods and
2 equally-entitled agents having valuation function $\vhat$.
Then the allocation $A = (\{1, 2\}, \{3\})$ is PROP but not M1S.
\end{lemma}
\begin{proof}
$A$ is PROP since $\vhat(A_1) = 4+4\eps$, $\vhat(A_2) = 3+2\eps$, and the PROP share is $2+3\eps$.
Suppose $A$ is M1S and agent 2's M1S certificate for $A$ is $B$.
$B_2 \neq \emptyset$, since $B$ is EF1-fair to agent 2.
Moreover, $\vhat(B_2) \le \vhat(A_2) = 3 + 2\eps$, so $B_2 = \{3\}$. Hence, $B_1 = \{1, 2\}$.
However, $\min_{g \in B_1} \vhat(B_1 \setminus \{g\}) = 4 + 4\eps > 3 + 2\eps = \vhat(B_2)$.
Hence, agent 2 is not EF1-satisfied by $B$, which contradicts the fact that $B$
is agent 2's M1S certificate for $A$. Hence, $A$ is not M1S.
\end{proof}

\begin{example}[PROP1 $\nfimplies$ PROPm]
\label{cex:ud:prop1-not-propm}
Consider a unit-demand function $v$ over 2 goods such that $v(1) = 30$ and $v(2) = 10$.
Let $0 \le \eps < 1$ and $\vhat(X) \defeq v(X) + |X|\eps$ for all $X \subseteq [2]$.
Consider a fair division instance with 2 goods and
2 equally-entitled agents having valuation function $\vhat$.
Then the allocation $A \defeq (\{1, 2\}, \emptyset)$ is PROP1 but not PROPm.
\end{example}

\begin{example}[PROP $\nfimplies$ PPROP]
\label{cex:ud:prop-not-pprop}
Consider a unit-demand function $v$ over 3 goods such that $v(1) = 30$ and $v(2) = v(3) = 11$.
Let $0 \le \eps < 1$ and $\vhat(X) \defeq v(X) + |X|\eps$ for all $X \subseteq [3]$.
Consider a fair division instance with 3 goods and
3 equally-entitled agents having valuation function $\vhat$.
Then the allocation $A \defeq (\{1\}, \{2\}, \{3\})$ is PROP but no alloction is PPROP.
\end{example}

\begin{example}[MMS $\nfimplies$ EF1, PROPm]
\label{cex:ud:mms-not-ef1-propm}
Consider a unit-demand function $v$ over 5 goods such that
$v(1) = 200$ and $v(2) = 30$ and $v(3) = v(4) = v(5) = 10$.
Let $0 \le \eps < 1$ and $\vhat(X) \defeq v(X) + |X|\eps$ for all $X \subseteq [5]$.
Consider a fair division instance with 5 goods and
4 equally-entitled agents having valuation function $\vhat$.
Then the allocation $A \defeq (\{1, 2\}, \{3\}, \{4\}, \{5\})$
is MMS but not EF1 and not PROPm.
(The maximin share is $10 + \eps$. The proportional share is $50 + (5/4)\eps$.)
\end{example}

\begin{example}[PROPm $\nfimplies$ PROPavg]
\label{cex:ud:propm-not-propavg}
Consider a unit-demand function $v$ over 3 goods such that
$v(1) = 75$, $v(2) = 30$, and $v(3) = 10$.
Let $0 \le \eps < 1$ and $\vhat(X) \defeq v(X) + |X|\eps$ for all $X \subseteq [3]$.
Consider a fair division instance with 3 goods and
3 equally-entitled agents having valuation function $\vhat$.
Then the allocation $A \defeq (\{1, 3\}, \{2\}, \emptyset)$ is PROPm but not PROPavg.
(The proportional share is $25+\eps$.)
\end{example}

\subsection{Submodular Valuations}
\label{sec:cex-extra:submod}

\begin{definition}[PMRF]
\label{defn:pmrf}
Let $P = (P_1, \ldots, P_k)$ be a partition of a finite set $M$.
Define $\PMRF(P)$ as the function $f: 2^M \to [|M|]$ where
$f(X) \defeq \sum_{i=1}^k \boolone(X \cap P_i \neq \emptyset)$.
(Here $\boolone(C)$ is 1 if condition $C$ is true and 0 otherwise.)
For any $\eps \in \mathbb{R}_{\ge 0}$, define $\PMRF_{+\eps}(P)$
as the function $f$ where $f(X) = \PMRF(P)(X) + |X|\eps$.
\end{definition}

\begin{lemma}
\label{thm:pmrf-submod}
For any partition $P$ over a finite set $M$,
$f \defeq \PMRF_{+\eps}(P)$ is submodular, and $f(g \mid X) \in \{\eps, 1+\eps\}$
for any $X \subseteq M$ and $g \in M \setminus X$.
\end{lemma}
\begin{proof}[Proof sketch]
$\PMRF(P)$ is the rank function of a partition matroid, and hence, is submodular.
% TODO: prove.
The submodularity of $f$ follows from \cref{thm:ud-perturb}.
\end{proof}

\begin{example}[EF $\nfimplies$ MMS]
\label{cex:ef-not-mms-pmrf}
Let $P \defeq (\{r_1, r_2\}, \{g_1, g_2\})$, $\eps \in \mathbb{R}_{\ge 0}$,
and $v \defeq \PMRF_{+\eps}(P)$.
Consider a fair division instance with 2 equally-entitled agents
having identical valuation function $v$.
Then the maximin share is at least $2 + 2\eps$,
as exemplified by the partition $(\{r_1, g_1\}, \{r_2, g_2\})$.
The allocation $P$ is EF but not MMS,
since each agent's value for her own bundle is $1 + 2\eps$.
\end{example}

\begin{example}[EF1 $\nfimplies$ MXS]
\label{cex:ef1-not-mxs-pmrf}
Let $P \defeq (\{r_1, r_2\}, \{g_1, g_2\}, \{b\})$ and $v \defeq \PMRF(P)$.
Consider a fair division instance with 2 equally-entitled agents
having identical valuation function $v$.
Then the allocation $A \defeq (\{r_1, r_2\}, \{g_1, g_2, b\})$ is EF1
but agent 1 is not MXS-satisfied by $A$.
\end{example}

\begin{example}[MEFS $\nfimplies$ EF1]
\label{cex:mefs-not-ef1-pmrf}
Let $P \defeq (\{a\}, \{b\}, \{c\}, \{d\}, \{e_1, e_2\}, \{f_1, f_2\}, \{g_1, g_2\})$.
Let $v \defeq \PMRF_{+\eps}(P)$ for $0 \le \eps \le 1/2$.
Consider a fair division instance with 2 equally-entitled agents
having identical valuation function $v$.
Let $A \defeq (\{a, e_1, f_1, g_1\}, \{b, c, d, e_2, f_2, g_2\})$.
Since $v(A_2) = 6 + 6\eps$ and $v(A_1) = 4 + 4\eps$,
agent 2 is envy-free in $A$ but agent 1 is not EF1-satisfied by $A$.
However, agent 1 is MEFS-satisfied by $A$, and her MEFS-certificate for $A$ is
$B \defeq (\{a, b, c, d\}, \{e_1, e_2, f_1, f_2, g_1, g_2\})$,
since $v(B_1) = v(A_1) = 4 + 4\eps$ and $v(B_2) = 3 + 6\eps$.
\end{example}

\begin{example}
\label{cex:eef-not-ef1-pprop-pmrf}
Let $M \defeq \{a_j\}_{j=1}^6 \cup \{b_j\}_{j=1}^6 \cup \{c_j\}_{j=1}^4$.
Let $P$ be a partition of $M$ where $X \in P$ iff $X = \{a_j, b_j\}$ for some $j \in [6]$
or $X = \{c_j\}$ for some $j \in [4]$.
Let $v \defeq \PMRF_{+\eps}(P)$ for $0 \le \eps < 1/6$.
Consider a fair division instance with 3 equally-entitled agents
having identical valuation function $v$.

Let $A \defeq (\{a_j\}_{j=1}^6, \{b_j\}_{j=1}^6, \{c_j\}_{j=1}^4)$.
Then $v(A_1) = v(A_2) = 6(1+\eps)$ and $v(A_3) = 4(1+\eps)$.
Then $A$ is not EF1 and not PPROP.
But $A$ is EEF since agents 1 and 2 are envy-free,
and agent 3's EEF-certificate is $B$ where
$B_1 \defeq \{a_1, a_2, a_3, b_1, b_2, b_3\}$,
$B_2 \defeq \{a_4, a_5, a_6, b_4, b_5, b_6\}$, and $B_3 = A_3$.
Then $v(B_3) = 4(1+\eps)$ and $v(B_1) = v(B_2) = 3 + 6\eps < 4$.
\end{example}

\begin{example}[PROP $\nfimplies$ M1S]
\label{cex:prop-not-m1s-uniform-matroid}
Let $v: 2^{[10]} \to \mathbb{R}_{\ge 0}$ be a function defined as $v(X) \defeq \min(6, |X|)$.
Then $v$ is submodular (since $v$ is the rank function of a uniform matroid)
and $v$ has binary marginals (i.e., $v(g \mid X) \in \{0, 1\}$
for all $X \subseteq [10]$ and $g \in [m] \setminus X$).
Let $0 \le \eps < 1$ and $\vhat(X) \defeq v(X) + |X|\eps$.

Consider a fair division instance with 2 equally-entitled agents
having identical valuation function $\vhat$.
Then $A \defeq ([4], [10] \setminus [4])$ is PROP but agent 1 is not M1S-satisfied.
\end{example}

\begin{lemma}[MMS $\nfimplies$ APS for $n=2$, Remark 2 of \cite{babaioff2023fair}]
\label{cex:mms-not-aps-n2-submod}
Let $C \defeq \{\{1, 2, 3\}, \{1, 5, 6\}, \{2, 4, 6\}, \{3, 4, 5\}\}$.
Let $v: 2^{[6]} \to \mathbb{R}_{\ge 0}$ be a function defined as
\[ v(X) \defeq \begin{cases}
0 & \text{ if } X = \emptyset
\\ 2 & \text{ if } |X| = 1
\\ 4 & \text{ if } |X| = 2
\\ 5 & \text{ if } |X| = 3 \text{ and } X \not\in C
\\ 6 & \text{ if } X \in C \text{ or } |X| \ge 4
\end{cases}. \]
Let $0 \le \eps < 1/3$ and define $\vhat$ as $\vhat(X) \defeq v(X) + |X|\eps$.
Then $\vhat$ is submodular.

Consider a fair division instance with 2 equally-entitled agents
having identical valuation function $\vhat$.
The maximin share is $5 + 3\eps$, but the APS is $6 + 3\eps$.
Moreover, there exists an MMS allocation but no allocation is APS.
\end{lemma}
\begin{proof}
For any $X \subseteq [6]$, and $g \in [6] \setminus X$, we have
\begin{tightenum}
\item $|X| \le 1 \implies \vhat(g \mid X) = 2 + \eps$.
\item $|X| = 2 \implies \vhat(g \mid X) \in \{1 + \eps, 2 + \eps\}$.
\item $|X| = 3 \implies \vhat(g \mid X) \in \{\eps, 1 + \eps\}$.
\item $|X| = 4 \implies \vhat(g \mid X) = \eps$.
\end{tightenum}
Hence, $\vhat(g \mid X)$ decreases with $|X|$, so $\vhat$ is submodular.

The MMS is less than 6 because for every set in $C$,
its complement is not in $C$.
The APS is $6 + 3\eps$ by the dual definition of APS (\cref{defn:aps-dual}),
by setting $x_S = 1/4$ when $S \in C$.

For identical valuations, an MMS allocation always exists.
For identical valuations, in any allocation, some agent gets a bundle of value at most the MMS.
Since the APS is strictly greater than the MMS, that agent is not APS-satisfied.
Hence, no allocation is APS.
\end{proof}

\subsection{Subadditive Valuations}
\label{sec:cex-extra:subadd}

\begin{lemma}
\label{thm:floor-ceil-value}
For any $k \in \mathbb{N}$, define the set functions $f_1, f_2: 2^{[m]} \to \mathbb{R}$ as
$f_1(S) = \floor{|S|/k}$ and $f_2(S) = \ceil{|S|/k}$, respectively.
Then $f_1$ is superadditive, $f_2$ is subadditive,
and for any $S \subseteq [m]$, $g \in [m] \setminus S$, and $i \in [2]$,
we have $f_i(g \mid S) \in \{0, 1\}$.
\end{lemma}

\begin{example}
\label{cex:gaps-not-ef1-prop1-subadd}
Let $([2], [9], (v_i)_{i=1}^2, \eqEnt)$ be a fair division instance where
$v_i(S) \defeq \ceil{|S|/4}$ for all $i \in [2]$ and $S \subseteq [9]$.
Then $v_i$ is subaditive for $i \in [2]$ by \cref{thm:floor-ceil-value}.
Let $A$ be an allocation where agent 1 has 3 goods and agent 2 has 6 goods.
Then $A$ is GAPS+EFX+PROPx but not EF1 or PROP1.

Each agent's APS is $\le 1$, because if we price each good at $1/9$,
then a bundle is affordable iff its value is at most 1.
Each agent's proportional share is $3/2$.
Note that $A$ is not \EFXZero{} (c.f.~\cref{defn:efx0-goods});
the difference between EFX and \EFXZero{} is crucial here.
\end{example}

\begin{example}
\label{cex:ef1-not-prop1-subadd}
Let $([2], [7], (v_i)_{i=1}^2, \eqEnt)$ be a fair division instance where
$v_i(S) \defeq \ceil{|S|/4}$ for all $i \in [2]$ and $S \subseteq [7]$.
Then $v_i$ is subaditive for $i \in [2]$ by \cref{thm:floor-ceil-value}.
Let $A$ be an allocation where agent 1 has 2 goods and agent 2 has 5 goods.
Then $A$ is GAPS+EF1 but not PROP1.
(Note that in our definition of PROP1 (c.f.~\cref{defn:prop1}), we require
$v_i(A_i \cup \{g\}) > v_i([m])/2$ for some $g \in [m] \setminus A_i$,
not just $v_i(A_i \cup \{g\}) \ge v_i([m])/2$.
This nuance is crucial to this counterexample.)
\end{example}

\begin{lemma}[optimal bin packing is subadditive]
\label{thm:bin-packing}
Let $(s_j)_{j=1}^m$ be a sequence where $0 \le s_j \le 1$ for all $j \in [m]$
($s_j$ is called the \emph{size} of item $j$).
For any $X \subseteq [m]$, define $\optBP_s(X)$ as the smallest integer $k$
such that a partition $(P_1, \ldots, P_k)$ of $X$ exists
such that $\sum_{j \in P_i} s_j \le 1$ for all $i \in [k]$.
Then $\optBP_s$ is subadditive and has binary marginals,
i.e., $\optBP_s(X \cup \{g\}) - \optBP_s(X) \in \{0, 1\}$.
\end{lemma}
\begin{proof}
Pick any two disjoint sets $X, Y \subseteq [m]$.
Let $k_X \defeq \optBP_s(X)$ and $k_Y \defeq \optBP_s(Y)$.
Let $(P_1, \ldots, P_{k_X})$ and $(Q_1, \ldots, Q_{k_Y})$ be the corresponding partitions.
Then $(P_1, \ldots, P_{k_X}, Q_1, \ldots, Q_{k_Y})$ is a feasible partition of $X \cup Y$,
so $\optBP_s(X \cup Y) \le k_X + k_Y$.
\end{proof}

\begin{example}
\label{cex:gaps-not-prop1-subadd}
Let $s = (1/5, 1/5, 3/5, 3/5, 3/5, 3/5)$. Let $0 \le \eps < 1/3$.
Consider a fair division instance with 3 equally-entitled agents
having an identical valuation function $v$ over 6 goods.
$v(X) \defeq \optBP_s(X) + |X|\eps$
($\optBP_s$ is defined in \cref{thm:bin-packing}).
Then by \cref{thm:bin-packing}, $v$ is subadditive,
and for every $X \subseteq [6]$ and $g \in [6] \setminus X$,
we have $v(g \mid X) \in \{\eps, 1+\eps\}$.

Then allocation $A = (\{1, 2\}, \{3, 4\}, \{5, 6\})$ is EFX, pairwise PROP, and groupwise APS
(set the price of the first two goods to $1/2$ and the last four goods to 1).
However, agent 1 is not PROP1-satisfied in $A$, since $v([6]) = 4 + 6\eps$,
and $v(A_1 \cup \{g\}) = 1 + 3\eps$ for all $g \in [6] \setminus A_1$

(Note that unlike \cref{cex:ef1-not-prop1-subadd}, this example
works even for the weaker definition of PROP1.)
\end{example}

\begin{lemma}[PAPS+PPROP $\nfimplies$ PROPm or M1S for binary subadd]
\label{cex:paps-pprop-not-propm-binary-subadd}
Let $A_1$, $A_2$, and $A_3$ be pairwise-disjoint sets where
$A_1$ contains 7 goods of size $0.6$, $A_2$ contains 21 goods of size $0.4$,
and $A_3$ contains 21 goods of size $0.4$. Let $M \defeq A_1 \cup A_2 \cup A_3$.
Let $v(X) \defeq \optBP(X)$ for all $X \subseteq M$.
Consider a fair division instance with 3 equally-entitled agents
having the same valuation function $v$ over $M$.
Then allocation $A \defeq (A_1, A_2, A_3)$ is EFX+PPROP+PAPS,
but $A$ is not PROPm-fair or M1S-fair to agent 1.
\end{lemma}
\begin{proof}
First, we show that $v(S) = \max(|S \cap A_1|, \ceil{|S|/2})$ for any $S \subseteq M$.
To see this, suppose $S$ contains $a$ items of size $0.6$ and $b$ items of size $0.4$.
If $b \le a$, we require $a$ to pack the items and $\ceil{(a+b)/2} \le a$.
If $b \ge a$, then we require $\ceil{(a+b)/2}$ bins to pack the items and $a \le (a+b)/2$.
As a corollary, we get that for all $S \subseteq M \setminus A_1$,
we have $v(S \mid A_1) = \max(0, \ceil{(|S|-|A_1|)/2})$.

$v(A_2) = v(A_3) = 11$ and $v(A_1) = 7$, so agents 2 and 3 are envy-free.
For any $S \subseteq A_2$, $v(S \mid A_1) > 0 \iff |S| \ge 8$,
and $|S| \ge 8 \implies |A_2 \setminus S| \le 13 \implies v(A_2 \setminus S) \le 7 = v(A_1)$.
Hence, agent 1 doesn't EFX-envy agent 2. Similarly, she doesn't EFX-envy agent 3.
Hence, $A$ is EFX.

$v(A_2 \cup A_3) = 21$, and $v(A_1 \cup A_2) = v(A_1 \cup A_3) = 14$.
Hence, $A$ is pairwise PROP. We will now show that $A$ is pairwise APS.
Set the price of each item to 1.
It's easy to see that the allocation restricted to agents 2 and 3 is APS.
Now suppose we restrict the allocation to agents 1 and 2.
Then $S \subseteq A_1 \cup A_2$ is affordable iff $|S| \le |A_1 \cup A_2|/2 = 14$.
For $|S| = 14$, we get $v(S) = 7$, so the APS is at most 7.
Hence, $A$ is pairwise APS satisfied.

$v(M)/3 = 25/3 = 8 + 1/3$.
For $S \subseteq A_2$ such that $|S| = 8$, we get that $v(S \mid A_1) = 1$,
but $v(A_1 \cup S) = 8 < v(M)/3$. Hence, $A$ is not PROPm-fair to agent 1.
Also, $v(A_1 \cup \{g\}) = v(A_1) = 7$ for all $g \in M \setminus A_1$,
so $A$ is also not PROP1-fair to $A$.

Suppose $A$ is M1S-fair to agent 1. Let $B$ be her M1S-certificate for $A$.
Then $v(B_1) \le v(A_1) = 7$ and $B$ is EF1-fair to agent 1.
$v(B_1) \le 7$ implies $|B_1| \le 14$. Since $|M| = 49$, we get that
$\max(|B_2|, |B_3|) \ge \ceil{|M - B_1|/2} \ge 18$.
Hence, some agent $i \in \{2, 3\}$ has a bundle of value at least $9$.
If one good is removed from their bundle, its value can drop by at most $8$,
so $B$ is not EF1-fair to agent 1, which is a contradiction.
Hence, $A$ is not M1S-fair to agent 1.
\end{proof}

\begin{lemma}[GMMS $\nfimplies$ APS]
\label{cex:gmms-not-aps-binary-subadd}
Let there be 15 items of sizes
$65/96$, $31/96$, $31/96$, $31/96$, $23/96$, $23/96$, $23/96$, $17/96$,
$11/96$, $7/96$, $7/96$, $7/96$, $5/96$, $5/96$, $5/96$.
Consider a fair division instance with 3 equally-entitled agents having
the same valuation function $v \defeq \optBP_s$ over the 15 items (c.f.~\cref{thm:bin-packing}).
Then $v$ is subadditive and has binary marginals.
Also, the AnyPrice share is at least 2, but the maximin share is 1.
Hence, no APS allocation exists, but the leximin allocation is GMMS.
\end{lemma}
\begin{proof}
The total size of the items is $3 \times 97/96$.
By Lemma C.1 of \cite{babaioff2023fair}, no $3$-partition of the items exists
such that the total size of each partition is at least $97/96$.
Hence, in every 3-partition $M$, some bundle has total size at most 1, so that bundle has value at most 1.
So, the maximin share is 1.

By Lemma C.1 of \cite{babaioff2023fair}, there exist 6 sets $(S_j)_{j=1}^6$ of items
such that each set has total size $97/96$ (and hence value 2) and each item appears in exactly two sets.
Hence, the APS is at least 2 (c.f.~\cref{defn:aps-dual}).
\end{proof}
